\begin{named}{Proposition~\ref{prop:research}(d), the segment at the cap}
Let \((F_r)_{r\in I}\) be a path on a compact interval \(I\) of strengths, along which \(\DT(r)\) and \(\tau(r)\) are
continuous and \(1/2<\tau(r)<1\), and let an indifferent investor mix over its best orders. Recruiting research at
the cap is an equilibrium at \(r\in I\) exactly when \(\DT(r)>2k\), \(\lambda_R(\DT(r))\ge\lb\) and
\(\tau(r)\le M_{\lb}\). The strengths with \(\DT(r)>2k\) and \(\lambda_R(\DT(r))\ge\lb\) are those of \(I\) at or
above a strength \(r_{\lb}\), and the strengths with \(\tau(r)\le M_{\lb}\) are those of \(I\) at or below a strength
\(r_C\). Recruiting research at the cap is therefore an equilibrium exactly on \([r_{\lb},r_C]\): the segment at the
cap starts where research reaches the cap and ends at the last strength at which the top posterior at the cap reaches
the preparation threshold.
\end{named}

\begin{proof}
By Proposition~\ref{prop:research}(a), recruiting research at \(r\) exists exactly when \(\DT(r)>2k\) and
\(\tau(r)\le M_\lambda\) at \(\lambda=\min\{\lambda_R(\DT(r)),\lb\}\), and this minimum is its research. Research is at
the cap exactly when \(\lambda_R(\DT(r))\ge\lb\), and the condition then reads \(\tau(r)\le M_{\lb}\). This is the
first claim.

By Proposition~\ref{prop:research}(b), \(\lambda_R\) is continuous and rises strictly on \(\DT>2k\), from \(0\) as
\(\DT\) falls to \(2k\). So there is a unique spread \(\DT^{\lb}>2k\) with \(\lambda_R(\DT^{\lb})=\lb\), and
\(\DT>2k\) with \(\lambda_R(\DT)\ge\lb\) holds exactly when \(\DT\ge\DT^{\lb}\). Along the path, \(\DT(r)\) and
\(\tau(r)\) do not fall (Lemma~\ref{lem:twonumbers}). Hence \(\{r\in I:\DT(r)\ge\DT^{\lb}\}\) contains every strength of
\(I\) above each of its members, and \(\{r\in I:\tau(r)\le M_{\lb}\}\) contains every strength of \(I\) below each of
its members. Both sets are closed, because \(\DT\) and \(\tau\) are continuous and \(I\) is compact. The first is
therefore \(\{r\in I:r\ge r_{\lb}\}\), with \(r_{\lb}\) its smallest member, and the second is
\(\{r\in I:r\le r_C\}\), with \(r_C\) its largest member. Their intersection is \([r_{\lb},r_C]\).

Before the cap, the posterior that a purchase must reach the threshold with is \(M_\lambda\) at
\(\lambda=\lambda_R(\DT(r))\), which moves with \(r\). From \(r_{\lb}\) on it is the fixed posterior \(M_{\lb}\), and
because \(\tau\) does not fall, \(\tau(r)\le M_{\lb}\) at some \(r\ge r_{\lb}\) gives it at every strength in
\([r_{\lb},r]\). The segment ends at the largest strength with \(\tau(r)\le M_{\lb}\). Where \(\tau\) is constant at the
value \(M_{\lb}\) on a stretch of strengths, that stretch belongs to the segment, and the segment ends at its right end.
\end{proof}

\paragraph{The segment at the cap in the example.}
Take the example at the prior \(1/2\), with \(a=\reg{res_a}\) and \(\lb=\reg{res_cap}\), on \(I=[r_0,r_2]\). Here
\(\DT(r)=(r-\ell)^2/(2r)\) rises strictly, and \(\tau(r)=(12r-3/4)/(20r-r^2-1)\) has
\(\tau'(r)=3(8r^2-r+2)/[2(20r-r^2-1)^2]>0\), so \(\tau\) rises strictly; also \(1/2<\tau(r)\le\tau(r_2)<1\). Since
\(M_{\lb}-\tau(r_1)=\reg{margin_res_tau_r1_B}\) and \(\tau(r_2)-M_{\lb}=\reg{margin_res_tau_r2_B}\), the last strength
with \(\tau(r)\le M_{\lb}\) is the unique root \(r_C\in(r_1,r_2)\) of \(\tau(r)=M_{\lb}=\eu/(1+\eu)\), the entry ceiling.
Because \((1-\rho)k/(2a\rho)>1\), the term \((\rho\DT m_\lambda-k)^+\) vanishes at the research root, and
\(\lambda_R(\DT(r))\ge\lb\) is \(\DT(r)\,m_{\lb}-k\ge2a\lb\). With \(m_{\lb}=1/(1+\eu)\), this is
\(\DT(r)\ge(2a\lb+k)(1+\eu)\), that is \(r\ge r_{\lb}=\fr\bigl((2a\lb+k)(1+\eu)\bigr)\), and
\(r_{\lb}<r_N<r_1<r_C\), with \(r_N-r_{\lb}=\reg{margin_res_rN_B}\). Recruiting research at the cap is therefore an
equilibrium exactly on \([r_{\lb},r_C]=[\reg{res_rcap},\reg{res_rC}]\).
